\documentclass[10pt,a4paper]{amsart}
\usepackage[margin=19mm]{geometry}
\usepackage{amsmath,amssymb,mathtools}
\usepackage[scr=rsfs,cal=euler]{mathalfa}
\usepackage{xfrac}
\usepackage[unicode,hidelinks,bookmarks=false]{hyperref}
\newcommand{\ie}{i.e.\ }
\newcommand{\cf}{cf.\ }
\newcommand{\resp}{respectively\ }
\newcommand{\Subsection}{\subsection}
\providecommand{\nobreakdash}{\nobreak\hbox{-}}
\setlength{\emergencystretch}{2em}
\multlinegap0pt
\usepackage{amssymb}
\usepackage{xcolor}
\usepackage{tikz}
\usepackage{algorithm}\usepackage{algpseudocode}
\newtheorem{lemma}{Lemma}
\newtheorem{theorem}{Theorem}
\title{A Min-Max Relation on Dicuts and Dijoins in Weighted Chordal Digraphs}
\author{G\'erard Cornu\'ejols, Siyue Liu,, R. Ravi}
\date{Source: arXiv:2501.10918v2 (CC BY 4.0)}
\begin{document}
\maketitle

\noindent\emph{Source and licence.} A Min-Max Relation on Dicuts and Dijoins in Weighted Chordal Digraphs. Version arXiv:2501.10918v2 (CC BY 4.0), distributed by arXiv under the Creative Commons Attribution 4.0 International licence, http://creativecommons.org/licenses/by/4.0/, as recorded in the arXiv licence metadata for this exact version and shown on the abstract page https://arxiv.org/abs/2501.10918v2. Full licence notice: \texttt{licenses/arxiv-2501.10918-CC-BY-4.0.txt}.\par\smallskip

\noindent\textit{Editorial scope.} Counted unit: the article's theorem thm:map\_back (source0.tex lines 284--286). The article's complete original proof of it is delivered with it (source0.tex lines 288--316, one \texttt{proof} environment of 29 lines). No mathematics is written, repaired or re-derived: the statement and the proof are the article's own lines, in the article's own notation, and no retained line is substituted or re-flowed.  Retained uncounted as the source-local context the proof needs: lines 197--199; lines 318--323.  One declared adaptation: exactly 2 ancillary strings inside the retained spans are omitted (image-inclusion commands and, where the source repeats a label, the repeated label definitions) and no other line differs from the source; the figure environments, with their placement, captions and labels, are kept, and the package records the omitted strings in fidelity receipts bound to the affected bodies.  No retained line contains a citation, so the wrapper carries no bibliography and no bibliography entry.  The retained lines contain the article's own diagram code, which the wrapper typesets with the standard tikz packages; no image file is delivered and the package declares no asset dependency.  Editorial formatting only, in addition: an AMS \texttt{amsart} preamble that restates the article's own declarations and macros used by the retained lines, namely the environments \texttt{lemma}, \texttt{theorem} on separate counters, together with the standard packages those lines need.  The retained environments print in this document as Lemma 1, Theorem 1 in order of appearance, whereas the article prints the same items with the numbering of its own full document.  The complete native source of the article as distributed in the arXiv e-print for this exact version is bundled unchanged as \texttt{sources/171789/source0.tex}.
\begin{lemma}\label{lemma:dicut_preserve}
    Every dicut in $D'$ is a dicut in $D$ restricted to the arcs in $A'$.
\end{lemma}

\begin{figure}[htbp]
	\centering
	
  \caption{In the unweighted case, for each color class, one dijoin in $\mathcal{J}'$ using the solid arc is mapped to one dijoin using the dashed arc in $\mathcal{J}$ of the same color.}
  \label{fig:weight_transfer2}
\end{figure}

    \begin{theorem}\label{thm:map_back}
        If there exists a packing of $\tau$ dijoins in $(D',w')$, then there exists a packing of $\tau$ dijoins in $(D,w)$.
    \end{theorem}

    \begin{proof}
    Let $\mathcal{J}'=\{J'_1,...,J'_\tau\}$ be a packing of $\tau$ dijoins in $(D',w')$. We map each one to a dijoin of $D$ according to the following two cases.
    
    \paragraph{\textbf{Case 1}:} $\delta(v)$ is not a dicut (see Figure \ref{fig:weight_transfer1}).
    
    By Lemma \ref{lemma:dicut_preserve}, the family of dicuts in $D'$ is the same as the family of dicuts in $D$ restricted to $A'$. For each $i\in\{1,...,s\}$ and $j\in\{s+1,...,k\}$, suppose $y_{ij}$ dijoins in $\mathcal{J}'$ use $v_iv_j$. Let $\mathcal{J}'_{ij}\subseteq \mathcal{J}'$ be a family of dijoins with  $|\mathcal{J}'_{ij}|=\min\{x_{ij}, y_{ij}\}$ such that each $J'\in \mathcal{J}'_{ij}$ uses $v_iv_j$. Let $\mathcal{J}_{ij}:=\{J'-v_iv_j+v_iv+vv_j\mid J'\in \mathcal{J}'_{ij}\}$. Every $J\in \mathcal{J}_{ij}$ is a dijoin in $D$ because every dicut containing $v_iv_j$ also contains one of $v_iv$ and $vv_j$. Replace $\mathcal{J}'_{ij}$ with $\mathcal{J}_{ij}$ in $\mathcal{J}'$ and go to the next $(i,j)$. For each $j\in\{s+2,...,k\}$, suppose $z_j$ dijoins in $\mathcal{J}'$ use $v_{s+1}v_j$. Let $\mathcal{J}'_j\subseteq \mathcal{J}'$ be a family of dijoins with $|\mathcal{J}'_j|=\min\{z_j,u_j-u_j'\}$ such that each $J'\in \mathcal{J}'_j$ uses $v_{s+1}v_j$. Let $\mathcal{J}_j:=\{J'-v_{s+1}v_j+vv_j\mid J'\in \mathcal{J}'_j\}$. Every $J\in \mathcal{J}_j$ is a dijoin in $D$ because every dicut containing $v_{s+1}v_j$ also contains $vv_j$. Replace $\mathcal{J}'_j$ with $\mathcal{J}_j$ in $\mathcal{J}'$ and go to the next $j$.
    Let $\mathcal{J}$ be the new family of dijoins.
    % obtained from $\mathcal{J}'$ with each $\mathcal{J}'_{ij}$ replaced by $\mathcal{J}_{ij}$ and each $\mathcal{J}'_j$ replaced by $\mathcal{J}_j$. 
    By the way we construct $w'$,  $\mathcal{J}$ is indeed a valid packing in $(D,w)$.
\begin{figure}[htbp]
	\centering
	
  \caption{In the unweighted case, for each color class, one dijoin in $\mathcal{J}'$ using the solid arc is mapped to one dijoin using the dashed arc in $\mathcal{J}$ of the same color.}
  \label{fig:weight_transfer1}
\end{figure}

    \paragraph{\textbf{Case 2}:} $\delta(v)$ is a dicut (see Figure \ref{fig:weight_transfer2}). 
    
    It is further required that the dijoins cover $\delta(v)$. 
    % Note that in this case, $s=0$ (since we assumed $\sum_{i=1}^s u_i\leq \sum_{i=s+1}^k u_i$) and there is no $\mathcal{J}_{ij}$. 
    Note that there is no $\mathcal{J}_{ij}$ in this case.
    Construct $\mathcal{J}_{j}'$ and $\mathcal{J}_{j}$ for $j\in\{2,...,k\}$ in the same way as in the previous case. Let $\mathcal{J}'_1:=\mathcal{J}'\backslash(\bigcup_{2\leq j\leq k} \mathcal{J}'_j)$
    be the remaining dijoins. To construct $\mathcal{J}$, we replace each $\mathcal{J}'_j$ by $\mathcal{J}_j$, $j\in\{2,...,k\}$. For each remaining dijoin $J'\in \mathcal{J}'_1$, we add one arc $vv_j$ to it for some $j\in \{1,...,k\}$. It suffices to argue that the number of remaining dijoins does not exceed the remaining total capacity of arcs incident to $v$.
    The crucial observation is that we can assume w.l.o.g. that every dijoin in $\mathcal{J}'$ uses at most one arc in $\{v_1v_j\mid j=2,...,k\}$. Indeed, if a dijoin $J'$ uses both $v_1v_i$ and $v_1v_j$ for some $2\leq i<j\leq k$, then $J'-v_1v_i$ is also a dijoin of $D'$ because every dicut containing $v_1v_i$ also contains $v_1v_j$. This implies that $\mathcal{J}_j'$'s, $j\in\{2,...,k\}$ are pairwise disjoint.
    Thus, $|\mathcal{J}_1'|=\tau-\sum_{j=2}^k |\mathcal{J}'_j|=\tau-\sum_{j=2}^k \min\{z_j, u_j\}\leq u_1+\sum_{j=2}^k (u_j-\min\{z_j, u_j\})$. The inequality follows from the fact that $\delta(v)$ is a dicut and thus has weight at least $\tau$. Since $u_1+\sum_{j=2}^k (u_j-\min\{z_j, u_j\})$ is the remaining total capacity of arcs incident to $v$, we have enough capacity for each remaining dijoin to include an arc incident to v. 
    
    % Let $\mathcal{J}_1:=\{J'+vv_1\mid J'\in \mathcal{J}'_1\}$. Let $\mathcal{J}$ be the family of dijoins obtained from $\mathcal{J}'$ with each $\mathcal{J}'_j$ replaced by $\mathcal{J}_j$, $j\in\{1,...,k\}$. Clearly, each $J\in\mathcal{J}$ is a dijoin. It remains to argue that $\mathcal{J}$ is indeed a valid packing under weight $w$. In particular, we need to show $|\mathcal{J}_1|\leq u_1$. 
    % The crucial observation is that we can assume w.l.o.g. that every dijoin in $\mathcal{J}'$ uses at most one arc in $\{v_1v_j\mid j=2,...,k\}$. Indeed, if a dijoin $J'$ uses both $v_1v_i$ and $v_1v_j$ for some $2\leq i<j\leq k$, then $J'-v_1v_i$ is also a dijoin of $D'$ because every dicut containing $v_1v_i$ also contains $v_1v_j$. Thus, $|\mathcal{J}_1|=|\mathcal{J}|-\sum_{j=2}^k |\mathcal{J}_j|=\tau-\sum_{j=2}^k u_j\leq u_1$. The inequality follows from the fact that $\delta(v)$ is a dicut and thus has weight at least $\tau$. 
    \end{proof}

\end{document}
